% 条件付き独立性の弦図（Fritz, arXiv:1908.07021, Def. 12.1 / Def. 12.8）。
% 図は下から上に読む。余域の順は X ⊗ Z ⊗ Y で固定する。
\documentclass[tikz,border=10pt,11pt]{standalone}
\usepackage[match]{luatexja-fontspec}
\setmainjfont{Noto Serif CJK JP}
\setsansjfont{Noto Sans CJK JP}
\usetikzlibrary{calc}

\definecolor{ink}{HTML}{1C1917}
\definecolor{muted}{HTML}{57534E}
\definecolor{wirec}{HTML}{292524}
\definecolor{panelbg}{HTML}{F7F6F3}
\definecolor{panelline}{HTML}{E7E5E4}
\definecolor{sfill}{HTML}{FEF3C7}
\definecolor{sedge}{HTML}{B45309}
\definecolor{pfill}{HTML}{FFEDD5}
\definecolor{pedge}{HTML}{C2410C}
\definecolor{qfill}{HTML}{E0F2FE}
\definecolor{qedge}{HTML}{0369A1}
\definecolor{rule}{HTML}{D6D3D1}

\begin{document}
\begin{tikzpicture}[
  x=1cm, y=1cm,
  every node/.style={inner sep=1.5pt, text=ink},
  wire/.style={draw=wirec, line width=1.7pt, line cap=round, line join=round},
  edge/.style={draw=sedge, line width=1.15pt, line join=miter},
  kern/.style={
    draw=pedge, fill=pfill, line width=1.05pt, rounded corners=2pt,
    minimum width=1.58cm, minimum height=0.78cm,
    inner sep=2pt, outer sep=0pt, font=\sffamily
  },
  kernq/.style={kern, draw=qedge, fill=qfill},
  var/.style={inner sep=1pt, outer sep=2pt},
  title/.style={font=\large, inner sep=2pt},
  sub/.style={font=\small, text=muted, inner sep=1pt},
  tag/.style={font=\small, text=muted, inner sep=1pt},
  panel/.style={draw=panelline, fill=panelbg, rounded corners=5pt, line width=0.8pt}
]

  % ---------- パネル ----------
  \draw[panel] (0.15,3.72) rectangle (11.55,11.62);
  \draw[panel] (11.78,3.72) rectangle (20.05,11.62);
  \draw[panel] (0.15,-2.05) rectangle (11.55,3.52);
  \draw[panel] (11.78,-2.05) rectangle (20.05,3.52);

  \node[title] at (5.85,11.18) {$X \perp Y \mid Z$ の定義（copy--discard の等式）};
  \node[sub] at (5.85,10.58) {同時状態 $p:I\to X\otimes Z\otimes Y$};
  \draw[rule, line width=0.6pt] (0.45,10.28) -- (11.25,10.28);

  \node[title] at (15.92,11.18) {conditional\_product の構成};
  \node[sub] at (15.92,10.58) {$p_{XZ}$ と $q_{YZ}$（共通の周辺 $p_Z=q_Z$）};
  \draw[rule, line width=0.6pt] (12.08,10.28) -- (19.75,10.28);

  % ============================================================
  % 左パネル: p = p_Z ; copy_Z ; (p(x|z) ⊗ id_Z ⊗ p(y|z))
  % ============================================================

  % 左辺。三本の出力を持つ同時状態。
  \draw[edge, fill=sfill]
    (0.48,4.48) -- (4.02,4.48) -- (3.58,5.92) -- (0.92,5.92) -- cycle;
  \node at (2.25,5.05) {$p$};
  \node[tag] at (2.25,4.16) {同時状態};

  \node[var] (Lx) at (1.18,9.78) {$X$};
  \node[var] (Lz) at (2.25,9.78) {$Z$};
  \node[var] (Ly) at (3.32,9.78) {$Y$};
  \draw[wire] (Lx.south) -- (1.18,5.92);
  \draw[wire] (Lz.south) -- (2.25,5.92);
  \draw[wire] (Ly.south) -- (3.32,5.92);

  \node at (4.84,7.15) {\scalebox{1.55}{$=$}};

  % 右辺。周辺、三叉の copy、両側の条件付き核。
  \draw[edge, fill=sfill]
    (7.60,4.48) -- (9.14,4.48) -- (8.37,5.62) -- cycle;
  \node at (8.37,4.90) {$p_Z$};
  \node[tag] at (8.37,4.16) {周辺};

  \node[kern] (px) at (6.44,8.22) {$p(x{\mid}z)$};
  \node[kern] (py) at (10.30,8.22) {$p(y{\mid}z)$};
  \node[var] (Rx) at (6.44,9.78) {$X$};
  \node[var] (Rz) at (8.37,9.78) {$Z$};
  \node[var] (Ry) at (10.30,9.78) {$Y$};

  \coordinate (copy) at (8.37,6.22);
  \draw[wire] (8.37,5.62) -- (copy);
  \draw[wire] (copy) to[out=108,in=-90] (px.south);
  \draw[wire] (copy) to[out=72,in=-90] (py.south);
  \draw[wire] (Rz.south) -- (copy);
  \draw[wire] (Rx.south) -- (px.north);
  \draw[wire] (Ry.south) -- (py.north);
  \node[tag, anchor=west, fill=panelbg, inner sep=1.5pt] at (8.58,9.22) {$\mathrm{id}_Z$};
  \fill[wirec] (copy) circle (5.2pt);
  \node[tag, anchor=west] at (8.68,5.90) {copy};

  % ============================================================
  % 右パネル: conditional_product(p_XZ, q_YZ)
  % ============================================================
  \begin{scope}[xshift=7.55cm]
    \draw[edge, fill=sfill]
      (7.60,4.48) -- (9.14,4.48) -- (8.37,5.62) -- cycle;
    \node at (8.37,4.90) {$p_Z$};
    \node[tag] at (8.37,4.16) {共通の周辺};

    \node[kern] (cpx) at (6.44,8.22) {$p(x{\mid}z)$};
    \node[kernq] (cqy) at (10.30,8.22) {$q(y{\mid}z)$};
    \node[var] (Cx) at (6.44,9.78) {$X$};
    \node[var] (Cz) at (8.37,9.78) {$Z$};
    \node[var] (Cy) at (10.30,9.78) {$Y$};

    \coordinate (ccopy) at (8.37,6.22);
    \draw[wire] (8.37,5.62) -- (ccopy);
    \draw[wire] (ccopy) to[out=108,in=-90] (cpx.south);
    \draw[wire] (ccopy) to[out=72,in=-90] (cqy.south);
    \draw[wire] (Cz.south) -- (ccopy);
    \draw[wire] (Cx.south) -- (cpx.north);
    \draw[wire] (Cy.south) -- (cqy.north);
    \node[tag, anchor=west, fill=panelbg, inner sep=1.5pt] at (8.58,9.22) {$\mathrm{id}_Z$};
    \fill[wirec] (ccopy) circle (5.2pt);
    \node[tag, anchor=west] at (8.68,5.90) {copy};
  \end{scope}

  % ============================================================
  % 下段左: Z が自明なら独立は周辺の積
  % ============================================================
  \node[title, font=\normalsize] at (5.85,3.16)
    {$Z$ が自明なら $X \perp Y$ は周辺の積};

  % パネルを縦に伸ばしたので、図は見出しの下の空きに置く。
  \begin{scope}[yshift=-1.00cm]
    \draw[edge, fill=sfill]
      (1.55,0.58) -- (3.95,0.58) -- (3.62,1.62) -- (1.88,1.62) -- cycle;
    \node at (2.75,1.00) {$p$};
    \node[var] (sx) at (2.15,2.48) {$X$};
    \node[var] (sy) at (3.35,2.48) {$Y$};
    \draw[wire] (sx.south) -- (2.15,1.62);
    \draw[wire] (sy.south) -- (3.35,1.62);

    \node at (4.70,1.45) {\scalebox{1.35}{$=$}};

    \draw[edge, fill=sfill] (5.55,0.58) -- (6.85,0.58) -- (6.20,1.62) -- cycle;
    \node at (6.20,0.96) {$p_X$};
    \node[var] (mx) at (6.20,2.48) {$X$};
    \draw[wire] (mx.south) -- (6.20,1.62);

    \draw[edge, fill=sfill] (7.55,0.58) -- (8.85,0.58) -- (8.20,1.62) -- cycle;
    \node at (8.20,0.96) {$p_Y$};
    \node[var] (my) at (8.20,2.48) {$Y$};
    \draw[wire] (my.south) -- (8.20,1.62);

    \node[sub] at (10.00,1.35) {$p=p_X\otimes p_Y$};
  \end{scope}

  % ============================================================
  % 下段右: 読み方、記号、構成式
  % ============================================================
  \node[title, font=\normalsize] at (15.92,3.16) {読み方・記号・構成式};

  \node[sub, anchor=west] at (12.05,2.46)
    {図は下から上に読む（下が域 $I$、上が余域）。};

  % copy。記号列の上端は y=1.86 に揃える。
  \fill[wirec] (12.28,1.57) circle (4.4pt);
  \draw[wire] (12.28,1.34) -- (12.28,1.57);
  \draw[wire] (12.28,1.57) -- (12.05,1.84);
  \draw[wire] (12.28,1.57) -- (12.28,1.86);
  \draw[wire] (12.28,1.57) -- (12.51,1.84);
  \node[tag, anchor=west] at (12.62,1.57) {copy};

  % discard（接地）
  \begin{scope}[shift={(14.55,1.80)}]
    \draw[wire] (0,0.10) -- (0,-0.02);
    \draw[wire] (-0.20,-0.02) -- (0.20,-0.02);
    \draw[wire] (-0.13,-0.14) -- (0.13,-0.14);
    \draw[wire] (-0.07,-0.26) -- (0.07,-0.26);
  \end{scope}
  \node[tag, anchor=west] at (14.88,1.57) {discard};

  % 状態
  \draw[edge, fill=sfill] (16.55,1.40) -- (17.15,1.40) -- (16.85,1.80) -- cycle;
  \node[tag, anchor=west] at (17.28,1.57) {状態};

  % 核
  \draw[pedge, fill=pfill, line width=1.05pt, rounded corners=1.5pt]
    (18.15,1.40) rectangle (18.85,1.78);
  \node[tag, anchor=west] at (18.98,1.57) {核};

  % 構成式は折り返さず、2 行に分けて記号列とも注とも離す。
  \node[sub] at (15.92,0.78)
    {$\mathrm{conditional\_product}(p_{XZ},q_{YZ}) =$};
  \node[sub] at (15.92,-0.08)
    {$p_Z \mathbin{;} \mathrm{copy}_Z \mathbin{;}
      (p(x{\mid}z)\otimes\mathrm{id}_Z\otimes q(y{\mid}z))$};
  \node[tag] at (15.92,-0.88)
    {$p_Z$ は $X$ と $Y$ を discard した周辺};
  \node[tag] at (15.92,-1.58)
    {Fritz, arXiv:1908.07021};
\end{tikzpicture}
\end{document}
